\chapter{Gen4 TopoChain：规格与数学}
\label{ch:gen4}

\section{定位}

TopoChain 是\textbf{Gear-free} 第四代 CDC：XOR-LFSR stride $\wedge$ 同构 $\hat\beta_0$ $\wedge$ 可选 $\hat\beta_1$。  
与 Hom（\texttt{0x20000}）\textbf{独立 dedup 域}；切点公式\textbf{冻结}——Limit Push 只允许 bit-identical 加速。

\begin{table}[htbp]
\centering
\caption{Chain 标识}
\label{tab:chain-id}
\begin{tabular}{@{}ll@{}}
\toprule
项 & 值 \\
\midrule
运行时 & \envvar{EBBACKUP\_CDC\_ALGO=topochain} \\
feature & \texttt{kBackupFeatureTopoChain = 0x40000} \\
\texttt{topo\_variant} & 2 \\
\texttt{chunk\_algo\_id} & 2（与 Hom 同族 ID，由 variant 区分） \\
并行 & \envvar{EBBACKUP\_TOPOCHAIN\_PARALLEL\_SCAN=1} \\
SPEC & \filepath{TOPO\_CHAIN\_GEN4\_SPEC.md} \\
\bottomrule
\end{tabular}
\end{table}

\section{三层切点公式}

量化与 LFSR：
\begin{align}
Q(b) &= b \gg q,\qquad q\in\{0,\ldots,7\}\quad\text{（默认 }0\text{）}\\
S_t &= \mathrm{rotl}_1(S_{t-1})\oplus Q(b_t)\\
M &= 2^{L}-1,\qquad L=\texttt{chain\_stride\_log}\in[8,22]\\
\mathrm{strideOK}_t &= (S_t \mathbin{\&} M)=0.
\end{align}

拓扑层（与 Hom 同构）：
\begin{align}
k_i &= Q(b_i)\\
\mathrm{topoOK}_t &= \bigl(\Delta\hat\beta_0\bigr)_t \neq 0
\equiv \texttt{ChainPrimaryDelta}\neq 0.
\end{align}

环检测（Tier-2，InitRepo 默认开）：
\begin{align}
\hat\beta_{1,\mathrm{GF2}}
&= \max\bigl(0,\, E - V + C\bigr)
\quad\text{（无向 1-骨架，顶点=出现过的键）}\\
\mathrm{beta1OK} &= (\hat\beta_1>0)\lor(\neg\texttt{enable\_beta1})\\
\mathrm{cut}_t &= \mathrm{strideOK}_t\land\mathrm{topoOK}_t\land\mathrm{beta1OK}\land(t\ge\mathrm{pos}+\min).
\end{align}

\begin{figure}[htbp]
\centering
\begin{tikzpicture}[
  box/.style={rectangle, rounded corners=2pt, draw=Gen4Color, fill=Gen4Color!10,
    align=center, font=\footnotesize, minimum height=0.8cm, minimum width=1.65cm},
  arr/.style={-{Stealth}, thick}]
  \node[box] (b) {byte};
  \node[box, right=0.4cm of b] (q) {$Q$};
  \node[box, right=0.4cm of q] (s) {$S_t$};
  \node[box, right=0.4cm of s] (st) {stride};
  \node[box, right=0.4cm of st] (b0) {$\Delta\hat\beta_0$};
  \node[box, right=0.4cm of b0] (b1) {$\hat\beta_1$};
  \node[box, right=0.4cm of b1] (c) {cut};
  \draw[arr] (b)--(q)--(s)--(st)--(b0)--(b1)--(c);
\end{tikzpicture}
\caption{Tier-0 / 1 / 2 串行门控。}
\label{fig:chain-tiers}
\end{figure}

\section{与 Hom 的同构定理（草图）}

\begin{designbox}[同构草图]
Hom 键 $=G[b]\mathbin{\&}0xFF$；Chain 键 $=Q(b)$。固定窗 $w$ 上两者都构造路径 1-骨架边集，且
\texttt{ChainPrimaryDelta} 与 \texttt{PrimaryTopoDelta} \textbf{字面同一更新公式}。\\
故 $\Delta\hat\beta_0$ 事件逻辑同构；差异仅 Tier-0：Gear mask $\leftrightarrow$ XOR-LFSR stride。
\end{designbox}

期望 stride 间距 $\approx M+1=2^L$。标定选择 $L$ 使 mean chunk $\in[0.85,1.15]\times\mathrm{avg}$（含 beta1 过滤效应）。

\section{超块}

\begin{table}[htbp]
\centering
\caption{Chain 超块语义}
\label{tab:chain-sb}
\begin{tabular}{@{}llp{6.8cm}@{}}
\toprule
字段 & 域 & 用途 \\
\midrule
\texttt{topo\_table\_seed} & 随机 u32 & LFSR 初态（\textbf{非} Gear） \\
\texttt{topo\_reserved[0:1]} & $L\in[8,22]$ & stride\_log \\
\texttt{topo\_reserved[2]} & $q\le7$ & quant\_q \\
\texttt{topo\_reserved[3]} & bit0 & \texttt{kChainFeatureBeta1} \\
\bottomrule
\end{tabular}
\end{table}

MVP 仓无 beta1 bit：\textbf{禁止}就地升级，须新 InitRepo。

\section{并行 scan 语义}

段长 256\,KiB。串行预热 $S$ 至各段起点 $s_i$，多线程探测，自左向右 join：左段命中则 cancel 右段。\\
跨 feed carry 时 \texttt{force\_serial}。语义仍取\textbf{全局最早} cut —— 与串行 bit-identical。

\section{复杂度}

非 stride：$O(1)$ LFSR；stride：$O(w)$ 键重建 $+\,O(w)$ $\hat\beta_1$。\\
设 stride 概率 $p\approx 2^{-L}$，摊销
\begin{equation}
T/\mathrm{byte}\approx T_{\mathrm{LFSR}} + p\cdot O(w).
\end{equation}

\section{门禁}

\texttt{L5\_topochain\_ab}：\texttt{vs\_stream}$\ge0.97$，\texttt{scan\_ns/probe}$\le1.30$。
